% ID: SEL-D0017-BB-E03
% ESTADO: validada. S018. Revisión matemática, didáctica y visual.
\textbf{a) Empalme en cero.}\par\nopagebreak[4]
El valor y el límite izquierdo son $2$. El límite derecho es $-b$, así que la continuidad exige $b=-2$. Las derivadas de los tramos son $2x+a$ y
\[
\frac{(x-1)-(x+b)}{(x-1)^2}=-\frac{1+b}{(x-1)^2}.
\]
Con continuidad, las pendientes laterales en cero son $a$ y $-(1+b)=1$. Por tanto, $\boxed{a=1,\ b=-2}$.

\textbf{b) Tangente con $a=1$, $b=-2$.}\par\nopagebreak[4]
Como $f(0)=2$ y $f'(0)=1$, la ecuación punto-pendiente da $\boxed{y=x+2}$.

\textbf{c) Asíntotas con $a=1$, $b=1$.}\par\nopagebreak[4]
En el tramo positivo se cumple $(x+1)/(x-1)=1+2/(x-1)$. Los límites en $1$ son $-\infty$ por la izquierda y $+\infty$ por la derecha. El límite hacia $+\infty$ vale $1$. Las asíntotas son
\[
\boxed{x=1},\qquad \boxed{y=1\quad(x\to+\infty)}.
\]
En $0$ los límites son finitos. Hacia $-\infty$, el tramo cuadrático crece en valor absoluto más deprisa que cualquier recta, lo que descarta una asíntota rectilínea en esa dirección.
